% Draft manuscript text derived from:
% 04_reconcile_manual_event_catalogs.ipynb
% S03_validate_catalogue_with_infrapy_bartlett.ipynb
% S03_validate_catalogue_with_infrapy_bartlett_continuous.ipynb
%
% The BibTeX key `infrapy` is provisional and should be replaced by the
% key used for the archived InfraPy software release in mybibfile.bib.

\subsection{Manual arrival picking and reduced-time event association}
\label{sec:methods_catalogue_association}

Infrasonic arrivals were picked manually on the three pressure channels using
Antelope \texttt{dbpick}.  Catalogue reconstruction retained non-deleted
phase-\(N\) picks within the accident interval.  The archived channel names
HD1--HD3 were mapped to the corresponding current StationXML channel names
DD1--DD3 without altering the archived pick times.  This procedure retained
384 picks: 66 on DD1, 162 on DD2, and 156 on DD3.  Unix epoch seconds were used
throughout to preserve the original millisecond timing.

Before association, each pick was reduced to the DD2 reference time using the
surveyed source--receiver geometry and the meteorologically estimated effective
sound speed:

\begin{equation}
t_{i,\mathrm{red}} = t_i - \frac{r_i-r_{\mathrm{DD2}}}{c_{\mathrm{eff}}},
\label{eq:reduced_pick_time}
\end{equation}

where \(r_i\) is the SLC-40--sensor distance and
\(c_{\mathrm{eff}}=350.977\ \mathrm{m\,s^{-1}}\).  The resulting corrections
were 0.090107, 0, and 0.012867~s for DD1, DD2, and DD3, respectively.  These
corrections were used only for event association; the waveform sample times
were not changed.

Reduced picks were sorted chronologically.  Beginning with the earliest
unassigned pick, picks within a following association window were grouped as a
candidate cluster.  At most one pick per channel was retained, selecting the
pick closest to the cluster seed when necessary.  A cluster was accepted only
when it contained picks from at least two distinct infrasound channels, and its
reduced event time was defined as the median of the selected reduced pick
times.  Picks that did not satisfy the two-channel criterion remained
unassociated rather than being forced into an event.

Association windows from 0.008 to 0.060~s were tested.  The number of
two-or-more-channel candidates increased from 128 at 0.008~s to 153 at
0.030~s, and reached a stable plateau of 154 candidates between 0.040 and
0.060~s.  We adopted 0.040~s because it was the smallest value on this plateau.
The single candidate added between 0.030 and 0.040~s had coherent DD2--DD3
waveforms after the predicted shift (maximum correlation 0.879) and a robust
stack peak signal-to-noise ratio (SNR) of 8.0, supporting its retention.  Of
the 154 candidates, one DD1--DD2 candidate was excluded following targeted
waveform review because it lacked both a DD3 pick and a coherent DD3 waveform
counterpart (DD2--DD3 correlation 0.092).  The resulting detection catalogue
contained 153 events.  Event detection was kept separate from subsequent
waveform-quality classification; no amplitude, SNR, or beamforming threshold
was used to define the catalogue count.

\subsection{Waveform-quality measurements}
\label{sec:methods_waveform_quality}

Waveform quality was measured consistently for every retained event using the
baseline-corrected pressure records.  The three traces were aligned using the
surveyed differential distances and \(c_{\mathrm{eff}}\), independently of the
subsequent beamforming result, and combined using a sample-by-sample median.
The signal interval extended from \(-0.08\) to \(+0.25\)~s relative to the
reduced event time, and the noise interval from \(-0.50\) to \(-0.15\)~s.
Peak SNR was defined as the maximum absolute median-stack amplitude divided by
the robust noise scale

\begin{equation}
\sigma_{\mathrm{MAD}} = 1.4826\,
\mathrm{median}\left(\left|p-\mathrm{median}(p)\right|\right).
\end{equation}

An RMS SNR was also calculated as the ratio of signal-window to noise-window
root-mean-square amplitudes.  Finally, the maximum normalized correlation
between the fixed-geometry-aligned DD2 and DD3 traces was found while allowing
a residual lag of \(\pm0.020\)~s.  The peak SNR and DD2--DD3 correlation were
used to define descriptive waveform-quality subsets: a usable subset required
peak SNR \(\geq3\) and correlation \(\geq0.5\), whereas a core subset required
peak SNR \(\geq5\) and correlation \(\geq0.7\).  These are operational quality
divisions rather than event-detection criteria.  RMS SNR was retained as a
continuous diagnostic because it depends strongly on window duration and can
be modest for short, high-amplitude impulses.

\subsection{Independent Bartlett beamforming validation}
\label{sec:methods_infrapy_validation}

The catalogue and array measurements were independently evaluated using the
frequency-domain Bartlett beamformer implemented in LANL InfraPy version
0.4.1.0 \citep{infrapy}.  The three baseline-corrected pressure channels were
processed directly as an ObsPy stream using their surveyed coordinates.  The
analysis band was 2--20~Hz and the beam window length was 0.50~s.  For the
event-centred analysis, one window was centred on each reduced catalogue time.
Back azimuth was searched from \(170^{\circ}\) to \(230^{\circ}\) in
\(0.5^{\circ}\) increments, and trace velocity from 280 to
440~\(\mathrm{m\,s^{-1}}\) in 2~\(\mathrm{m\,s^{-1}}\) increments.  The
surveyed back azimuth to SLC-40 was \(199.438^{\circ}\); the meteorological
effective speed of 350.98~\(\mathrm{m\,s^{-1}}\) was retained as an external
reference and was not imposed on the beam search.

At each grid point, InfraPy returned normalized Bartlett coherence \(C\).  For
three channels, the reported Fisher statistic was calculated as

\begin{equation}
F = (N-1)\frac{C}{1-C}, \qquad N=3.
\label{eq:bartlett_f_statistic}
\end{equation}

An empirical screening threshold was obtained from 153 evenly spaced,
non-overlapping 0.50-s windows preceding the first catalogue event.  The
99.5th percentile of this background distribution was \(F=3.023\).

To test catalogue recovery without using the catalogue times to position the
analysis windows, the complete common three-channel interval was also scanned
with 0.50-s windows stepped by 0.10~s.  For computational efficiency, this
blind scan used a \(1^{\circ}\) back-azimuth grid and a
5~\(\mathrm{m\,s^{-1}}\) trace-velocity grid over the same search limits.
Blind candidates were defined as local maxima above the empirical \(F\)
threshold, separated by at least 0.10~s and having a prominence of at least
twice the robust background \(F\)-statistic scale.  Only after peak detection
were blind candidates matched greedily and one-to-one to catalogue events
within \(\pm0.25\)~s.  InfraPy's sequential detection grouping was also
retained to describe broad episodes of coherent activity, but these episodes
were not interpreted as individual explosions because overlapping windows can
merge closely spaced events.

\subsection{Catalogue construction and waveform quality}
\label{sec:results_catalogue_quality}

The final catalogue comprises 153 events between 13:07:15.927 and
13:33:54.396~UTC, a span of 1598.47~s (26.64~min).  Ninety events were supported
by DD2 and DD3, and 63 by all three infrasound sensors; thus every retained
event included picks on both DD2 and DD3.  After travel-time
reduction, the median span among picks assigned to an event was 0.00487~s and
the 95th percentile was 0.0173~s.

The median robust peak SNR was 7.86.  Of the 153 events, 146 had peak SNR
\(\geq3\), 120 had peak SNR \(\geq5\), and 63 had peak SNR \(\geq10\).  The
median RMS SNR was 2.23.  Median DD2--DD3 correlation was 0.876; 137 events had
correlation \(\geq0.5\), 117 had correlation \(\geq0.7\), and 70 had
correlation \(\geq0.9\).  Applying the joint descriptive criteria yielded 132
events in the usable subset and 100 in the core subset.  For 151 of 153 events,
the best residual DD2--DD3 lag after the fixed geometric correction was no
larger than one sample (0.004~s) in magnitude.

\subsection{Event-centred Bartlett validation}
\label{sec:results_event_centred_bartlett}

The event-centred InfraPy analysis yielded a median \(F\) statistic of 5.04,
with 113 of 153 event windows exceeding the empirical pre-event threshold of
3.023.  The median back azimuth was \(200.675^{\circ}\), compared with the
surveyed SLC-40 direction of \(199.438^{\circ}\), and the median absolute
back-azimuth error was \(1.43^{\circ}\).  Among all events, 145 were within
\(5^{\circ}\) and 151 were within \(10^{\circ}\) of SLC-40.  Among the 113
events above the empirical \(F\) threshold, 112 were within \(5^{\circ}\) and
all were within \(10^{\circ}\).

The median trace velocity was 354.6~\(\mathrm{m\,s^{-1}}\), close to but not
identical with the meteorological effective speed of
350.98~\(\mathrm{m\,s^{-1}}\).  Trace-velocity scatter decreased strongly with
beam quality.  The 5th--95th percentile interval was 308--428~\(\mathrm{m\,s^{-1}}\)
for all events, 327--394~\(\mathrm{m\,s^{-1}}\) for events above the empirical
threshold, and 346--372~\(\mathrm{m\,s^{-1}}\) for the 42 events with
\(F\geq10\).  No event with \(F\geq5\) reached a trace-velocity grid boundary,
indicating that the extreme speeds were confined to weak, poorly constrained
windows rather than strong solutions truncated by the search limits.

\subsection{Blind continuous Bartlett validation}
\label{sec:results_continuous_bartlett}

The blind scan evaluated 17,839 overlapping windows and identified 418 local
\(F\)-statistic maxima above the deliberately permissive empirical threshold.
Greedy one-to-one matching within \(\pm0.25\)~s recovered 129 of the 153
catalogue events (84.3\%).  If sharing of a broad peak by closely spaced events
was allowed, 136 catalogue events had a blind peak within 0.25~s.  The 129
matched blind peaks had median \(F=8.08\), median back azimuth
\(200.623^{\circ}\), and median trace velocity
356.0~\(\mathrm{m\,s^{-1}}\); 126 were within \(5^{\circ}\) and all were within
\(10^{\circ}\) of SLC-40.

The 24 events not recovered by one-to-one matching were systematically weaker
than the recovered events.  Their median event-centred \(F\), peak SNR, RMS
SNR, and DD2--DD3 correlation were 2.19, 5.70, 1.59, and 0.606, respectively,
compared with 5.76, 8.85, 2.42, and 0.893 for recovered events.  Some nominally
missed events occurred in dense sequences and lay close to a broad blind peak
that had already been assigned to a neighboring catalogue event.

The 289 unmatched local maxima were predominantly marginal or occurred in the
coda of catalogue events.  At \(F\geq10\), 57 of 62 blind peaks matched the
catalogue, and the five unmatched peaks all occurred within 1~s of a catalogue
event.  Every one of the 31 peaks with \(F\geq20\) matched the catalogue.  No
unmatched peak more than 1~s from a catalogue event exceeded \(F=8.41\).
During the 183.6~s of the blind scan preceding the first catalogue event, 11
low-level maxima exceeded the permissive \(F=3.023\) threshold, but none
reached \(F=4\); their directions and trace velocities were not stable.

\subsection{Catalogue robustness and limitations of the independent validation}
\label{sec:discussion_catalogue_validation}

The reduced-time reconstruction supports a 153-event \emph{detection}
catalogue, while the waveform metrics describe a nested continuum of event
quality.  This distinction is important: a weak event need not satisfy a
high-SNR or three-channel beamforming criterion to be credible when it has
geometrically consistent manual picks on DD2 and DD3; waveform correlation and
beam coherence provide additional continuous measures of confidence rather
than binary definitions of event existence.  Conversely, DD1 provides useful
supporting information but should not
veto an event because it produced substantially fewer accepted picks and was
often less consistent than the other two sensors.

The InfraPy analysis provides independent support for the source direction and
for the moderate-to-strong portion of the catalogue.  Both event-centred and
blind results cluster tightly around the surveyed direction to SLC-40, whereas
the trace-velocity distribution narrows only as coherence increases.  The
broad speeds obtained for weak windows should therefore be interpreted as poor
resolution of the back-azimuth--slowness surface, not as evidence for similarly
large physical variations in atmospheric sound speed.  The approximately
355--357~\(\mathrm{m\,s^{-1}}\) median obtained from the better constrained
solutions is a trace velocity and should not be conflated with the independently
estimated 350.98~\(\mathrm{m\,s^{-1}}\) effective propagation speed.

The blind scan does not provide an independent proof that the exact event count
must be 153.  Its 0.50-s overlapping windows cannot resolve every event in the
densest parts of the sequence, while ringing, coda energy, and adjacent window
maxima can produce several blind peaks around one physical event.  This explains
both the unmatched weak catalogue events and the large raw count of 418 blind
maxima.  Nevertheless, the absence of any strong unmatched peak away from the
catalogue, together with complete matching of all peaks with \(F\geq20\),
provides no evidence for a missing moderate or strong explosion.  The blind
scan also found no coherent moderate or strong precursor during its 3.06-min
pre-event interval; this statement is limited to that interval and does not by
itself establish the absence of precursors throughout the longer archived
record.
